% TOP_PROBLEM: {"id": 3019, "title": "Kirby Problem 5.12", "category_id": 7, "category": {"id": 7, "name": "topology", "display_name": "Topology", "description": "Properties preserved under continuous deformations.", "slug": "topology", "order_index": 7, "created_at": "2026-07-31T15:26:25.671Z"}, "status": "open", "problem_number": "KP-5.12", "set_id": 11, "statusQualification": "Uses the usual finite-polyhedron interpretation of the Zeeman problem; compatible triangulations and subdivision are allowed for a polyhedral collapse."}
% FIRST_POSTED: 2026-10-01
% ENTRY_KIND: statement_only
% UnsolvedMath ID 3019; source catalogue code KP-5.12
% !TeX program = lualatex
% Definition and sources only; this file is not an LLM solution attempt.
\documentclass[11pt,a4paper]{article}
\usepackage[margin=25mm]{geometry}
\usepackage{fontspec}
\setmainfont{lmroman10-regular.otf}[BoldFont=lmroman10-bold.otf,
 ItalicFont=lmroman10-italic.otf,BoldItalicFont=lmroman10-bolditalic.otf]
\usepackage{amsmath,amssymb,mathtools}
\usepackage{unicode-math}
\setmathfont{latinmodern-math.otf}
\AtBeginDocument{\let\setminus\smallsetminus}
\usepackage{xurl,hyperref}
\hypersetup{colorlinks=true,urlcolor=blue}
\setcounter{secnumdepth}{0}
\setlength{\parindent}{0pt}
\setlength{\parskip}{5pt}
\setlength{\emergencystretch}{3em}
\begin{document}
\section*{Kirby Problem 5.12}\label{3019}
{\small UnsolvedMath ID 3019 · KP-5.12 · Topology · Kirby's Problems in Low-Dimensional Topology}
\subsection{Definitions and mathematical statement}
Let $K$ be a finite $2$-dimensional polyhedron, given by a
finite simplicial complex, and suppose $K$ is contractible:
its identity map is homotopic to a constant map. Write
$I=[0,1]$ and give the product $K\times I$ its natural
polyhedral structure, which admits finite triangulations of
dimension at most three.

An elementary simplicial collapse removes a simplex $\sigma$
together with a codimension-one face $\tau$ which is a free
face: $\sigma$ is the only maximal simplex containing
$\tau$. The remaining simplices form a subcomplex. A complex
is collapsible if a finite sequence of elementary collapses
reduces it to one vertex. For a polyhedron, the assertion
uses a compatible finite triangulation, allowing subdivisions.

Does $K\times I$ always collapse to a point in this sense?
This is the Zeeman conjecture. The product with an interval
adds one dimension to the available collapse process; no
further product factor or additional cell attachment is
permitted in the requested conclusion.

Each elementary collapse is a particularly controlled deformation
retraction, so collapsibility implies contractibility. The hypothesis
only supplies contractibility of $K$, and that weaker property
alone does not specify a collapse sequence. The conclusion
asks for a combinatorial simplification of the product itself.
The question ranges over all finite contractible $2$-complexes,
without assuming they embed in a $3$-manifold or arise as a
spine of one.
\subsection{Short English statement}
Does the product of every finite contractible two-dimensional polyhedron with an interval admit a collapse to a point?
\subsection{Sources}
\begin{itemize}
\item[S1] ULAM.AI and the UnsolvedMath Contributors, supplied problems.json, record 3019 (KP-5.12), Kirby's Problems in Low-Dimensional Topology. Catalogue identity and source transcription; CC BY 4.0.
\url{https://huggingface.co/datasets/ulamai/UnsolvedMath}.
\item[S2] K3: A New Problem List in Low-Dimensional Topology, Problem 5.12. Expanded formulation of the supplied catalogue statement.
\url{https://math.berkeley.edu/sites/default/files/surv-295-ruberman-watermarked-author-pdf.pdf}.
\end{itemize}
\end{document}
